\section{Preliminaries}
\label{sec:prelim}
%%AZ.8.29: might want to expand on the existing failure modes of BFMs here. I know it's in the workshop paper, but this paper should be mostly self-contained.
%%AZ.9.5: I would split this section up. The first part is just preliminaries/background, but the last two paragraphs are more important since they define the problem you are trying to solve. You could even fold them into the beginning of Section 2 (the section numbering is weird right now since your intro isn't under a section). 

\paragraph{Offline reward-free setting.}
We consider a controlled Markov process $\mathcal M=(\mathcal S,\mathcal A, \mu, P)$ and a dataset $\mathcal D=\{(s_i,a_i,s_i')\}_{i=1}^{N}$ collected without downstream reward labels.
Here, $\mathcal S$ denotes the state space, $\mathcal A$ denotes the action space, $\mu\in\Delta(\mathcal S)$ denotes the initial state distribution, and $P(\cdot\mid s,a)\in\Delta(\mathcal S)$ denotes the environment transition kernel.
Let $\rho$ denote the marginal distribution of next states induced by the offline data-collection process, from which the observed next states $\{s_i'\}_{i=1}^N$ are sampled.
Downstream tasks are specified by state rewards $r:\mathcal S\rightarrow\mathbb R$.

\paragraph{Successor measures.}
For a policy $\pi$ and discount factor $0<\gamma<1$, the normalized SM is the discounted distribution of future states:
\begin{equation}
    M^\pi(s,a,ds^+) :=(1-\gamma)\sum_{t\geq0}\gamma^t \mathbb P_\pi\!\left( s_{t+1}\in ds^+ \mid s_0=s,\,a_0=a \right),
    \label{eq:sm}
\end{equation}
Each term is the distribution of the state $t+1$ steps ahead after taking $a$ and following $\pi$, $s^+$ denotes a generic future state and $ds^+$ indicates the integration variable in $M^\pi(s,a,ds^+)$ and $\rho(ds^+)$.
Since $(1-\gamma)\sum_{t\geq0}\gamma^t=1$, $M^\pi(s,a,\cdot)$ is a probability measure.
It satisfies the Bellman recursion
\(
    M^\pi(s,a,\cdot) =(1-\gamma)P(\cdot\mid s,a) +\gamma\mathbb E_{\substack{ s'\sim P(\cdot\mid s,a)\\ a'\sim\pi(\cdot\mid s') }}\!\left[M^\pi(s',a',\cdot)\right],
    \label{eq:sm-bellman}
\)
and for any bounded state reward $r$, integrating the reward under the SM recovers the action value:
\begin{equation}
    Q_r^\pi(s,a) =\frac{1}{1-\gamma} \E_{s^+\sim M^\pi(s,a,\cdot)}[r(s^+)] =\frac{1}{1-\gamma} \int_{\mathcal S}r(s^+)\,M^\pi(s,a,ds^+).
    \label{eq:sm-value}
\end{equation}
\paragraph{Forward-Backward representations.}
FB~\citep{touati-neurips21a,touati-iclr23b} learns a family of policies $\pi_z$ indexed by task vectors $z\in\mathbb R^d$, together with a forward representation $F_\theta(s,a,z)\in\mathbb R^d$ and a backward representation $B_\theta(s^+)\in\mathbb R^d$, which factorize the density of the SM relative to the dataset state marginal $\rho$ as $ M_\theta(s,a,z,ds^+) := F_\theta(s,a,z)^\top B_\theta(s^+)\,\rho(ds^+).$
For a dataset transition $(s,a,s')$, an independently sampled state $s^+\sim\rho$, and $a'\sim\pi_z(\cdot\mid s')$, FB learns these representations from the Bellman equation through the squared regression objective
\begin{equation}
    \mathcal L_{\mathrm{FB}} = \E\Big[ \big( F_\theta(s,a,z)^\top B_\theta(s^+) - \gamma \bar F(s',a',z)^\top \bar B(s^+) \big)^2 - 2(1-\gamma)\, F_\theta(s,a,z)^\top B_\theta(s') \Big],
    \label{eq:fb-objective}
\end{equation}
where $\bar F$ and $\bar B$ denote target representations. 
For a downstream reward $r$, FB forms $z_r:=\E_{s\sim\rho}[r(s)B_\theta(s)]$. 
Substituting the factorization into \cref{eq:sm-value} gives the linear readout $Q_r^{\pi_z}(s,a)\approx F_\theta(s,a,z)^\top z_r/(1-\gamma)$

\paragraph{TD InfoNCE for policy evaluation.}
TD InfoNCE~\citep{zheng-iclr24a} estimates a policy's successor distribution using bootstrapped contrastive learning.
For a fixed policy $\pi$, let $h(s,a,\cdot)$ be a normalized density relative to $\rho$, and $\bar h$ a fixed target density.
With exact expectations and normalization, its evaluation objective is
\begin{equation}
        \mathcal L_{\mathrm{TD}}(h) :=-\E\Big[ (1-\gamma)\log h(s,a,s') +\gamma\E_{s^+\sim\rho}\!\left[ \bar h(s',a',s^+)\log h(s,a,s^+) \right] \Big],
        \label{eq:td-infonce-population}
\end{equation}
where $(s,a,s')\sim\mathcal D$ and $a'\sim\pi(\cdot\mid s')$.
The first term fits the observed next state; the second fits the target model's predicted continuation.
The practical objective approximates normalization and future-state expectations using sampled candidates.